\chapter{Método}

O primeiro objetivo é utilizar a ferramenta web scraper para extrair informações relevantes de diversas fontes da internet, como sites de notícias e redes sociais, a fim coletar dados relacionados a eventos e ao sentimento do mercado, de maneira a identificar tendências e padrões que influenciam o comportamento das ações. 

Uma vez coletadas as informações, o próximo objetivo é realizar o tratamento e a preparação adequada dos dados para servirem como entrada para o modelo CNN. Isso inclui a limpeza dos textos, a remoção de ruídos, a tokenização e a vetorização dos dados, garantindo assim que estejam em um formato adequado para a análise.

A sintetização desse modelo será conduzida em etapas sequenciais, de acordo com o processo de transformação informacional, de modo a garantir que as informações disponibilizadas pelo modelo possuam a adaptabilidade necessária para suprir os parâmetros exigidos pelos modelos de análise financeira.

O objetivo principal da aplicação da CNN é processar e classificar os dados de entrada transformados a fim de identificar padrões e características relevantes relacionadas ao contexto favorável ou desfavorável de cada ação. A saída do modelo CNN será uma pontuação ou probabilidade que indicará o grau de favorabilidade de cada ação analisada.

Após obter as pontuações ou probabilidades das ações, o objetivo é realizar o tratamento adequado desses dados de saída a fim de utilizá-los como parâmetros para um modelo final de alocação de ativos. Essa etapa envolverá a interpretação dos resultados da CNN e a definição de critérios para a seleção e alocação dos ativos na carteira de investimentos. Um exemplo de modelo final que pode ser utilizado é o modelo Black-Litterman, amplamente reconhecido na área financeira.

{\color{red}
\section{Comandos de Referências}


Comando $\backslash$cite:

\cite{NBR6024:2012}

\vspace{0.5cm}


Comando $\backslash$cite, com texto: 

\cite[p. 15]{NBR14724:2011}

\vspace{0.5cm}



Comando $\backslash$citeonline:

\citeonline{BORTOLUZZI}

\vspace{0.5cm}



Comando $\backslash$citeauthoronline:

\citeauthoronline{CONCEICAO}

\vspace{0.5cm}


Comando $\backslash$citeyear: 

\citeyear{GIL}

\vspace{0.5cm}



Uso do apud:

\apud{SIMOES}{GIL}

\vspace{0.5cm}


\cite{BU}

\vspace{0.5cm}

Para consultar a documentação de referências digite no terminal: \$ texdoc abntex2cite
}